\section{Mechanisation and regression evidence}\label{sec:mechanisation}

\subsection{Verified build}
The development is in Coq 8.18.0 \cite{coq818} and comprises roughly three and a half thousand lines of theory and nine hundred of examples. It builds with \texttt{coq\_makefile} and with \texttt{dune build}; \texttt{coqchk} succeeds on the whole project excluding one isolated file; and \texttt{Print Assumptions} reports each of the 120 Coq identifiers cited in this paper (92 of them in the ledger of \cref{app:ledger}) as closed under the global context. The one exception by design is the unrestricted functional converse of fibre factorisation, which needs classical choice and is confined to a single file that nothing else imports. Its constructive replacements are the relational form and, for finite domains with decidable equality, a direct construction. The main dependency closure is therefore axiom-free. \Cref{tab:trust} summarises the trust boundary.
\begin{table}[t]
\centering\small
\caption{Trust boundary and non-claims.}\label{tab:trust}
\begin{tabularx}{\textwidth}{@{}lY@{}}
\toprule
Checked mechanically & the theorems of this paper, in the Coq kernel; axiom-freeness of the main closure. \\
Trusted & the Coq kernel and extraction mechanism; the encoding of $\nat$ as OCaml integers in extracted code; the declared records themselves. \\
Tested, not proved & agreement of the extracted lineage audit with a handwritten audit (\cref{sec:regression}). \\
Not established & completeness of declared lineage; accuracy of raw inputs; adequacy of the state model; legitimacy or competence of an authority; decidability of general obligations. \\
\bottomrule
\end{tabularx}
\end{table}

\subsection{Legacy compatibility}
An earlier supplement (\cref{sec:seams}) is imported unchanged and built as an independent library. Its files match their recorded checksums, its own build-and-check procedure completes, and compatibility theorems recover its results from the generic calculus. Nothing in the earlier development was redefined.

\subsection{Extraction}\label{sec:extraction}
Certificate data live in $\Type$; their proofs live beside them. Extraction removes the proof fields and keeps the data (\cref{tab:survive}).
\begin{table}[t]
\centering\small
\caption{Evidence fields and whether they survive extraction.}\label{tab:survive}
\begin{tabularx}{\textwidth}{@{}lYc@{}}
\toprule
Component & Field & Survives \\
\midrule
Exhibition certificate & record id, identifiers, process, retained record and evaluator, coverage, custodian & yes \\
 & injectivity of identifiers; valuation factorisation & no \\
Lineage certificate & graph, coverage record with gaps, issued verdict & yes \\
 & well-formedness, L1--L3 & no \\
Construction certificate & seam inhabitant & yes \\
 & grounding equations & no \\
Obstruction & constructor tag; located data (seam element, node id) & yes \\
 & proofs (e.g.\ ancestor derivations) & no \\
Assessment & verdict, open reasons, gaps & yes \\
Maintenance and authority evidence & application-supplied data, via the summary interface & if in $\Type$ \\
\bottomrule
\end{tabularx}
\end{table}
We extracted the assessments of the examples to OCaml and wrote a consumer. For a certified assessment it reports the custodian (7), record identifier (1), lineage ground node (5) and audit verdict (\texttt{LineagePass}). For the copied model it reports a refutation at construction whose constructor names a lineage descent; for the maintenance model, a maintenance obstruction; for the model with both, both, in order; for missing and partial lineage records, an open assessment. The formal distinction thus survives compilation. It is materially stronger than attaching abstract propositions (\texttt{Exhibited}, \texttt{LineageGrounded}) to an ordinary relation, whose proofs would be erased.

Maintenance and authority evidence remain abstract types in the generic library. The generic calculus provides a \emph{summary interface}, a class assigning to each evidence type a function into a list of small extractable atoms (records, custodians, ground nodes, intervals, authorities, maintenance events, expiry, revalidation). Applications supply instances, so an extracted consumer can inspect more than the exhibition and lineage components without the kernel fixing an institutional schema. This is an observability extension, not a kernel correctness matter.

\subsection{Regression against a handwritten audit}\label{sec:regression}
The extracted checker (\cref{sec:lineage}) was compared with a handwritten OCaml lineage audit, by capturing the latter's printed output for the same records. This is \emph{differential regression testing}, not a proof of equivalence of the OCaml sources.

\paragraph{Case study.}
On five case-study records (live rows; backfill with a disclosed shared ancestor; backfill with it undisclosed; a copied seam valuation; a cyclic record) the two agree on all five, including the malformed cycle.

\paragraph{Unrestricted random records.}
On 3000 random records, including every malformation kind, they agree on all 3000. Only 89 of these were grounded, 1403 not grounded, 1508 malformed; the random population exercises the rejecting paths much more than the accepting path.

\paragraph{Stratified populations.}
To remedy this, we generate separate populations, each with a known \emph{intended} outcome. A graph passes only if (i) the extracted checker's verdict matches the intent, and (ii) its L1--L3 flags match the handwritten audit's output. The eleven \emph{shallow} populations, of 300 graphs each, are: accepted graphs (with random shared ancestry, dispositions and shuffled order); the same graphs with a coverage gap, intended $\mathsf{Open}$; one family for each single defect (descent from $d_A$, descent from $d_B$, undisclosed shared ancestor, no independent source); and one for each malformation (duplicate name, undeclared parent, distinguished node not derived, relevant input not raw, cycle). The handwritten audit has no notion of coverage, so the open population is checked against the intent and against the handwritten flags on the underlying graph.

\paragraph{Deep-ancestry populations.}
The mutation study below showed that the shallow populations are too shallow: their ancestor chains are short, so a saturation that stops early still finds every decisive ancestor. We therefore added six \emph{deep} populations of 300 graphs, in which the decisive ancestor is reachable only after several saturation rounds: the ground reaches its independent raw source through a chain of three to five derived nodes (depth at least four); descent from $d_A$ or $d_B$ is through a three-node branch (depth four); a shared derived ancestor is reached through chains of two to four nodes from each coordinate. The six populations are: accepted (L3 source at depth $\ge 4$), the same with a coverage gap ($\mathsf{Open}$), deep descent from $d_A$, deep descent from $d_B$, deep undisclosed shared ancestor, and deep no-independent-source. Together the seventeen populations comprise 5100 graphs, and all pass both criteria.

\paragraph{Mutation study.}
A test is only as good as its ability to fail. We textually mutated the \emph{extracted} checker in thirteen ways and recorded which suite detects each (\cref{tab:mutants}): the five-record case study (fixed), the unrestricted random population, the eleven shallow stratified populations, and the six deep populations.
\begin{table}[t]
\centering\small
\caption{Mutation study of the extracted lineage checker: $\times$ means the suite detects the mutant.}\label{tab:mutants}
\begin{tabular}{@{}lcccc@{}}
\toprule
Mutant & fixed & random & shallow & deep \\
\midrule
M1 drop the $d_B$ descent test (L1) & & $\times$ & $\times$ & $\times$ \\
M2 ignore dispositions (L2) & $\times$ & $\times$ & $\times$ & $\times$ \\
M3 raw shared ancestors count as undisclosed (L2) & $\times$ & $\times$ & $\times$ & $\times$ \\
M4 source need not be claim-relevant (L3) & & $\times$ & $\times$ & $\times$ \\
M5 source may be an ancestor of $d_A$ (L3) & & $\times$ & $\times$ & \\
M6 coverage gaps ignored & & & $\times$ & $\times$ \\
M7 skip relevance-is-raw check (WF) & & $\times$ & $\times$ & \\
M8 skip cycle detection (WF) & $\times$ & $\times$ & $\times$ & \\
M9 skip distinguished-derived check (WF) & & $\times$ & $\times$ & \\
M10 skip declared-parents check (WF) & & $\times$ & $\times$ & \\
M11 ancestors: parents only & $\times$ & $\times$ & $\times$ & $\times$ \\
M12 ancestors: one saturation step & $\times$ & $\times$ & & $\times$ \\
M13 ancestors: two saturation steps & & $\times$ & & $\times$ \\
\midrule
Mutants detected (of 13) & 5 & 12 & 11 & 8 \\
\bottomrule
\end{tabular}
\end{table}
All thirteen mutants are detected by the union of the suites, and by the stratified suite (shallow and deep together) alone. The individual suites are not equivalent, which is the argument for stratification. The random population detects twelve but misses M6, since the handwritten audit has no coverage; M6 is caught only by the open populations. The shallow populations detect eleven, and \emph{miss} the bounded-saturation mutants M12 and M13, because their graphs are shallow; the unrestricted random population detected those two because it happens to generate deeper graphs; and the deep populations, added in response, detect both. The fixed case study detects five and would not have found, for instance, the dropped $d_B$ descent test (M1).

The theorem of \cref{sec:lineage} already establishes checker correctness. These tests probe extraction, the encoding of records, the differential harness, and diagnostic agreement with the handwritten implementation; the mutation study measures whether that probe can fail.
